\documentclass{article}
\usepackage{bm}
\usepackage{mathtools}
\usepackage{amsthm}
\usepackage{fancyhdr}
\usepackage{float}
\usepackage{listings}
\lstset{breaklines=true}
\usepackage{graphicx}
\usepackage{xcolor}
\usepackage{titlesec}
\usepackage{titling}
\usepackage{tabularx}
\usepackage{enumitem}
\usepackage{amsmath}
\usepackage{amsfonts}
\usepackage{hyperref}
\usepackage[margin=1in]{geometry}

\hypersetup{
	colorlinks=true,
	linkcolor=blue,
	filecolor=magenta,
	urlcolor=blue,
	citecolor=blue,
	anchorcolor=blue
}

\renewcommand{\thesubsubsection}{\Roman{subsubsection}}
\title{SFWRENG 3O03 2023 Midterm Solutions}
\author{Matthew Farah, \href{mailto:farahm11@mcmaster.ca}{$\langle$farahm11@mcmaster.ca$\rangle$}, 400468251}
\date{\today}

\makeatletter
\setlist{listparindent=\parindent}
\newcolumntype{Y}{>{\centering\arraybackslash}X}

\pagestyle{fancy}
\fancyhead{}
\fancyhead[L]{Matthew Farah, $\langle$\href{mailto:farahm11@mcmaster.ca}{farahm11@mcmaster.ca}$\rangle$, 400468251}
\fancyhead[R]{SFWRENG 3O03 2023 Midterm Solutions}

\newcounter{problem}
\renewcommand{\theproblem}{\arabic{problem}}

\newenvironment{problem}[1][]{
	\ifx\\#1\\
		\refstepcounter{problem}
	\else
		\setcounter{problem}{#1}
	\fi
	\large{\par\noindent\textbf{Problem \theproblem}}\vspace{.5em}\par\noindent
}{\vspace{1.5em}}

\newenvironment{solution}{
	\vspace{0.4cm}\par\noindent\large\textbf{{Solution:}}\par\vspace{1em}
}{\vspace{2.5em}}

\newcounter{partslist}

\makeatletter
\newcommand{\parts@seriesname}{parts@\theproblem}

\newenvironment{parts}{
	\edef\parts@ser{\parts@seriesname}
	\ifcsname\parts@ser\endcsname
		\begin{enumerate}[resume=\parts@ser,label=\textbf{(\alph*)},leftmargin=3.3em,itemsep=0.4em,before=\large]
	\else
		\begin{enumerate}[series=\parts@ser,label=\textbf{(\alph*)},leftmargin=3.3em,itemsep=0.4em,before=\large]
		\expandafter\gdef\csname\parts@ser\endcsname{1}
	\fi
}{
	\end{enumerate}
}

\makeatother

\makeatletter

\newcounter{currentstep}[partslist]
\renewcommand{\thecurrentstep}{\Roman{currentstep}}

\newenvironment{steps}{
	\edef\steps@ser{steps@\theproblem}
	\ifcsname\steps@ser\endcsname
		\begin{enumerate}[
			resume=\steps@ser,
			label=\bfseries{(\Roman*)},
			leftmargin=3.15em,
			itemsep=0.4em,
			topsep=0.8em
		]
	\else
		\begin{enumerate}[
			series=\steps@ser,
			label=\bfseries{(\Roman*)},
			leftmargin=3.15em,
			itemsep=0.4em,
			topsep=0.8em
		]
		\expandafter\gdef\csname\steps@ser\endcsname{1}
	\fi
	\let\steps@saved@item\item
	\renewcommand{\item}{\stepcounter{currentstep}\steps@saved@item}
}{
	\let\item\steps@saved@item
	\end{enumerate}
}

\makeatother

\begin{document}

\maketitle
\pagestyle{fancy}

\begin{problem}
	We know that for a discrete system that $H(\omega) = e^{-i\omega} -e^{-2i\omega}$. What is the corresponding difference equation for this system?
\end{problem}

\begin{solution}
	We can use the fact $y(n) = H(\omega)e^{in\omega}$ when $x(n) = e^{in\omega}$ to recover the system's difference equation

	\begin{align*}
		H(\omega) &= e^{-i\omega} - e^{-2i\omega} \\
		H(\omega)e^{in\omega} &= e^{in\omega} \cdot \left( e^{-i\omega} - e^{-2i\omega} \right ) \\
		y(n) & = e^{-i\omega} \cdot e^{in\omega} - e^{-2i\omega} \cdot e^{in\omega} \\
		&= e^{i\omega(n - 1)} - e^{i\omega(n - 2)} \\
		&= x(n - 1) - x(n - 2)
	\end{align*}
\end{solution}

\begin{problem}
	What are the Fourier series coefficients of the signal $x(t) = 2 + \sin(t) + \cos(3t) + 2\cos(5t)$? You do not need to compute integrals to get full marks.
\end{problem}

%TODO: Consider not giving all coefficients newlines?

\begin{solution}
	We will rely on the fact that each term which makes up the signal is periodic, we can use the coefficients of the complex exponential expansions of sine and cosine ($sin(t) = \frac{e^{it} - e^{-it}}{2i}, \cos(t) = \frac{e^{it} - e^{-it}}{2})$ respectively) to find the Fourier series coefficients. Therefore:

	\begin{align*}
		&X_{-5} = 1 \\
		&X_{-4} = 0 \\
		&X_{-3} = \frac{1}{2} \\
		&X_{-2} = 0 \\
		&X_{-1} = \frac{-1}{2i} \\
		&X_0 = 2 \\
		&X_1 = \frac{i}{2i} \\
		&X_2 = 0 \\
		&X_3 = \frac{1}{2} \\
		&X_4 = 0 \\
		&X_5 = 1
	\end{align*}

	Notice here that:

	\begin{enumerate}
		\item The value of the Fourier series coefficient at index $j$ corresponds to the complex exponential expansion of the term in $x(n)$ with angular frequency $j$.
		\item Angular frequencies which do not appear in the signal are given a value of zero.
		\item $X_j = \bar{X}_{-j}$
	\end{enumerate}
\end{solution}

\begin{problem}
	Define the frequency response of a continuous system.
\end{problem}

\begin{solution}
	The frequency response, $H(\omega)$ of a continuous system $y(t)$ is defined as $y(t) = H(\omega)e^{it\omega}$ when $x(t) = e^{it\omega}$.
\end{solution}

\begin{problem}
	Use convolution to compute the output of a discrete system with impulse response $h(n)$ for the input $x(n) = e^{-in\omega}$. Note that the final result is the \emph{Discrete Time Fourier Series} (DTFT)
\end{problem}

\begin{solution}
	\begin{align*}
		y(n) &= \sum_{k = -\infty}^{\infty} \left( h(k)x(n - k) \right) \\
		&= \sum_{k = -\infty}^{\infty} \left( h(k)e^{i\omega(n - k)} \right) \\
		&= \sum_{k = -\infty}^{\infty} \left( h(k) \cdot e^{-in\omega} \cdot e^{ik\omega} \right) \implies \\
		y(n)e^{in\omega} &= \sum_{k = -\infty}^{\infty} \left( h(k) \cdot e^{ik\omega} \right) \implies \\
		H(\omega) &= \sum_{k = -\infty}^{\infty} \left( h(k) \cdot e^{ik\omega} \right)
	\end{align*}
\end{solution}

\begin{problem}
	Given the difference equation $y(n) = x(n) - y(n - 2)$, compute the system's output to the signal $x(n) = 1 + \cos(n\pi + \frac{\pi}{2}) + \sin(n\frac{\pi}{4}))$ using the its frequency response.
\end{problem}

\begin{solution}

	\begin{steps}
		\item Find the system's frequency response ($H(\omega)$).
	\end{steps}

	\begin{gather*}
		y(n) = x(n) - y(n - 2) \\\\
		H(\omega)e^{in\omega} = e^{in\omega} - H(\omega)e^{i\omega(n - 2))} \\\\
		H(\omega)e^{in\omega} = e^{in\omega} - H(\omega)e^{in\omega}e^{-2i\omega} \\\\
		H(\omega) = 1 - H(\omega)e^{-2i\omega} \\\\
		H(\omega) + H(\omega)e^{-2i\omega} = 1 \\\\
		H(\omega)(1 + e^{-2i\omega}) = 1 \\\\
		H(\omega) = \frac{1}{1 + e^{-2i\omega}}
	\end{gather*}

	\begin{steps}
		\item Compute $H(\omega), \, |H(\omega)|$, and $arg(H(\omega))$ for each $\omega$ in the signal $x(n)$.
	\end{steps}

	\begin{tabularx}{\linewidth}{Y | Y | Y | Y}
		$\omega$ & $H(\omega)$ & $|H(\omega)|$ & $arg(H(\omega))$ \\
		\hline
		0 & $1/2$ & $1/2$ & 0 \\
		$\pi$ & $1/2$ & $1/2$ & 0 \\
		$\pi/4$ & $1/{(1 - i)}$ & $\sqrt{2}/2$ & $-7\pi/4$
	\end{tabularx}

	\begin{steps}
		\item Use the formula $y(n) = |H(\omega)|\cos(n\omega + arg(H(\omega)))$ for each of the terms we computed.
	\end{steps}

	\begin{align*}
		\begin{split}
			y(n) &= |H(\omega_1)|\cos(n\omega_1 + arg(H(\omega_1))) + |H(\omega_2)|\cos(n\omega_2 + arg(H(\omega_2))) \\\\
			&+ |H(\omega_3)|\cos(n\omega_3 + arg(H(\omega_3))) \\\\
			&= \frac{1}{2}\cos(0n + 0) + \frac{1}{2}\cos(n\pi + 0) + \frac{\sqrt{2}}{2}\cos(\frac{pi}{4} + \frac{-7\pi}{4}) \\\\
			&= \frac{1}{2} + \frac{1}{2}\cos(\pi) + \frac{\sqrt{2}}{2}\cos(\frac{\pi}{4} + \frac{-7\pi}{4})
		\end{split}
	\end{align*}
\end{solution}

\begin{problem}
	Compute the impulse response and the frequency response of $y(n) = x(n - 3)$.

	\begin{steps}
		\item Compute the impulse response ($h(n)$).
	\end{steps}

	Substituting $x(n) = \delta(n)$, we find:

	\begin{align*}
		y(n) &= x(n - 3) \implies \\
		h(n) &= \delta(n - 3)
	\end{align*}

	\begin{steps}
		\item Compute the frequency response ($H(\omega)$).
	\end{steps}

	Using $y(n) = H(\omega)e^{in\omega}$ when $x(n) = e^{in\omega}$, we find:

	\begin{align*}
		y(n) &= x(n - 3) \\
		H(\omega)e^{in\omega} &= e^{i\omega(n - 3)} \\
		H(\omega)e^{in\omega} &= e^{in\omega} \cdot e^{-3i\omega} \\
		H(\omega) &= e^{-3i\omega}
	\end{align*}
\end{problem}

\begin{problem}
	What is the output y(n) given the system $y(n) = x(n) + {\alpha}y(n - 1)$ and input $x(n) = 1 \; \forall n \in \mathbb{N}$?
\end{problem}

\begin{solution}
	\begin{align*}
		y(0) &= 1 + \alpha(0) = 1 \\
		y(1) &= 1 + \alpha(1) = 1 + \alpha \\
		y(2) &= 1 + \alpha(1 + \alpha) = 1 + \alpha + {\alpha}^2 \\
		y(3) &= 1 + \alpha(1 + \alpha + {\alpha}^2) = 1 + \alpha + {\alpha}^2 + {\alpha}^3 \\
	\end{align*}

	Extrapolating, we can surmise that:

	\begin{gather*}
		y(n) = 1 + \alpha + {\alpha}^2 + \ldots + {\alpha}^n
	\end{gather*}
\end{solution}

\begin{problem}
	The Fourier series for the 4-periodic signal $x(n) = \sin(\frac{\pi}{2}n)$ was computed and yielded:
	\begin{gather*}
		X_0 = 0, \, X_1 = -i, \, X_2 = 0, \, X_3 = i \\
	\end{gather*}
	Compute the values of x(n) for one period.
\end{problem}

\begin{solution}
	\begin{steps}
		\item Compute $\omega_0$.
	\end{steps}

	Since we are given that $x(n)$ is 4-periodic, we can use the formula $\omega_0 = \frac{2\pi}{P}$ to find that:

	\begin{align*}
		\omega_0 &= \frac{2\pi}{4} \\
		&= \frac{\pi}{2}
	\end{align*}

	\begin{steps}
		\item Compute the values of $x(n)$ for one period.
	\end{steps}

	Using the formula:

	\begin{gather*}
		x(n) = \sum_{k = 0}^{p = 1}\left( X_k e^{ink\omega_0} \right) \\
	\end{gather*}

	we can compute the values of x(n) as follows:

	\begin{align*}
		x(0) &= \sum_{k = 0}^{3} \left( X_k e^{i(0)k\frac{\pi}{2}} \right) = (0 - i + 0 - i) = 0 \\
		x(1) &= \sum_{k = 0}^{3} \left( X_k e^{i(1)k\frac{\pi}{2}} \right) = (-ie^{\frac{\pi}{2}i} + ie^{\frac{3\pi}{2}i}) = 2 \\
		x(2) &= \sum_{k = 0}^{3} \left( X_k e^{i(2)k\frac{\pi}{2}} \right) = (-ie^{i\pi} + ie^{3i\pi}) = 0 \\
		x(3) &= \sum_{k = 0}^{3} \left( X_k e^{i(3)k\frac{\pi}{2}} \right) = (-ie^{\frac{3\pi}{2}i} + ie^{\frac{9\pi}{2}i}) = -2 \\
	\end{align*}
\end{solution}











\end{document}
